\clearpage
\subsection{Skalarprodukt in Koordinatendarstellung}
\label{subsec:vector_algebra_dot_product_coordinates}

Nun suchen wir noch eine Formel, um das Skalarprodukt zweier Vektoren,
die in ihrer Koordinatenform vorliegen zu berechnen. Diese Formel heißt
Koordinatenformel des Skalarprodukts.
Dazu stellen wir zunächst die zwei Vektoren $\vec a,\vec b\in\R^3$ mit
den kanonischen Einheitsvektoren dar:
\[
  \begin{aligned}
    \vec a
      &=\Cloze{10cm}{a_1\vec e_1+a_2\vec e_2+a_3\vec e_3},\\[1em]
    \vec b
      &=\Cloze{10cm}{b_1\vec e_1+b_2\vec e_2+b_3\vec e_3}.
  \end{aligned}
\]

Die kanonischen Einheitsvektoren sind orthonormal, d.h. sie sind 
paarweise senkrecht und haben den Betrag $1$. Daher gilt
\[
  \vec e_i\cdot\vec e_j
  =\begin{cases}
     1,&i=j,\\
     0,&i\ne j.
   \end{cases}
\]

Durch Einsetzen in das Skalarprodukt und anschliendes Ausmultiplizieren erhalten wir
\begin{ClozeArea}[34]{10}
  \[
    \begin{aligned}
      \vec a\cdot\vec b
      &=(a_1\vec e_1+a_2\vec e_2+a_3\vec e_3)
        \cdot(b_1\vec e_1+b_2\vec e_2+b_3\vec e_3)
        \vphantom{\overbrace{(\vec e_1\cdot\vec e_1)}^{=1}}\\[0.7em]
      &=a_1b_1\overbrace{(\vec e_1\cdot\vec e_1)}^{=1}
        +a_1b_2\overbrace{(\vec e_1\cdot\vec e_2)}^{=0}
        +a_1b_3\overbrace{(\vec e_1\cdot\vec e_3)}^{=0}\\[0.7em]
      &\quad
        +a_2b_1\overbrace{(\vec e_2\cdot\vec e_1)}^{=0}
        +a_2b_2\overbrace{(\vec e_2\cdot\vec e_2)}^{=1}
        +a_2b_3\overbrace{(\vec e_2\cdot\vec e_3)}^{=0}\\[0.7em]
      &\quad
        +a_3b_1\overbrace{(\vec e_3\cdot\vec e_1)}^{=0}
        +a_3b_2\overbrace{(\vec e_3\cdot\vec e_2)}^{=0}
        +a_3b_3\overbrace{(\vec e_3\cdot\vec e_3)}^{=1}\\[0.7em]
      &=a_1b_1+a_2b_2+a_3b_3
        \vphantom{\overbrace{(\vec e_1\cdot\vec e_1)}^{=1}}.
    \end{aligned}
  \]
\end{ClozeArea}

\begin{TheoremBox}{Koordinatenformel des Skalarprodukts}{vector_algebra_dot_product_coordinate_formula}
  Für das Skalarprodukt zweier Vektoren $\vec a,\vec b\in\R^n$ gilt
  \[
    \vec a\cdot\vec b
    =\begin{pmatrix}
      a_1\\a_2\\\vdots\\a_n
    \end{pmatrix}
    \cdot
    \begin{pmatrix}
      b_1\\b_2\\\vdots\\b_n
    \end{pmatrix}
    =a_1b_1+a_2b_2+\dots+a_nb_n
    =\sum_{i=1}^{n}a_i b_i.
  \]
  Komponenten mit gleichem Index werden multipliziert und anschließend werden
  die Produkte addiert.
\end{TheoremBox}

\begin{ExampleBox}{Skalarprodukt aus Koordinaten}{vector_algebra_coordinate_dot_product_example}
  \begin{enumerate}[label=\alph*),leftmargin=1.8em,itemsep=1em]
    \item
      \(\displaystyle
        \begin{pmatrix}1\\2\end{pmatrix}
        \cdot
        \begin{pmatrix}3\\4\end{pmatrix}
        =\Cloze{5cm}{1\cdot3+2\cdot4=11}.
      \)

    \item
      \(\displaystyle
        \begin{pmatrix}1\\2\\3\end{pmatrix}
        \cdot
        \begin{pmatrix}4\\5\\6\end{pmatrix}
        =\Cloze{7cm}{1\cdot4+2\cdot5+3\cdot6=32}.
      \)
  \end{enumerate}
\end{ExampleBox}
